\documentclass[10pt,a4paper]{amsart}
\usepackage[margin=19mm]{geometry}
\usepackage{amsmath,amssymb,mathtools}
\usepackage[scr=rsfs,cal=euler]{mathalfa}
\usepackage{xfrac}
\usepackage[unicode,hidelinks,bookmarks=false]{hyperref}
\newcommand{\ie}{i.e.\ }
\newcommand{\cf}{cf.\ }
\newcommand{\resp}{respectively\ }
\newcommand{\Subsection}{\subsection}
\providecommand{\nobreakdash}{\nobreak\hbox{-}}
\setlength{\emergencystretch}{2em}
\multlinegap0pt
\usepackage{amssymb}
\usepackage{xcolor}
\newtheorem{lemma}{Lemma}
\newtheorem{remark}{Remark}
\newtheorem{theorem}{Theorem}
\title{A Periodic Dichotomy in Linear Control Theory}
\author{Shichao Ye, Xingwu Zeng, Can Zhang}
\date{Source: arXiv:2604.04003v1 (CC BY 4.0)}
\begin{document}
\maketitle

\noindent\emph{Source and licence.} A Periodic Dichotomy in Linear Control Theory. Version arXiv:2604.04003v1 (CC BY 4.0), distributed by arXiv under the Creative Commons Attribution 4.0 International licence, http://creativecommons.org/licenses/by/4.0/, as recorded in the arXiv licence metadata for this exact version and shown on the abstract page https://arxiv.org/abs/2604.04003v1. Full licence notice: \texttt{licenses/arxiv-2604.04003-CC-BY-4.0.txt}.\par\smallskip

\noindent\textit{Editorial scope.} Counted unit: the article's theorem thm:analy\_riccati (source0.tex lines 1029--1035). The article's complete original proof of it is delivered with it (source0.tex lines 1045--1140, one \texttt{proof} environment of 96 lines, which the article places away from the statement). No mathematics is written, repaired or re-derived: the statement and the proof are the article's own lines, in the article's own notation, and no retained line is substituted or re-flowed.  Retained uncounted as the source-local context the proof needs: lines 172--186; lines 174--180; lines 188--201; lines 203--216; lines 222--264; lines 252--262; lines 992--1000.  Nothing was omitted from any retained span, so the package carries no omission record.  No retained line contains a citation, so the wrapper carries no bibliography and no bibliography entry.  No retained line contains a figure environment, an image-inclusion command or drawing code, so no image is delivered and the package declares no asset dependency.  Editorial formatting only, in addition: an AMS \texttt{amsart} preamble that restates the article's own declarations and macros used by the retained lines, namely the environments \texttt{lemma}, \texttt{remark}, \texttt{theorem} on separate counters, together with the standard packages those lines need.  The retained environments print in this document as Lemma 1, Remark 1, Theorem 1 in order of appearance, whereas the article prints the same items with the numbering of its own full document.  The complete native source of the article as distributed in the arXiv e-print for this exact version is bundled unchanged as \texttt{sources/197854/source0.tex}.
	\begin{lemma}\label{ric}
		The periodic matrix Riccati differential equation
		\begin{equation}\label{du12052}
			\begin{split}
				\dot{P}(t)&+A^*(t)P(t)+P(t)A(t)\\
				&-P(t)B(t)Q^{-1}(t)B^*(t)P(t)\\
				&+C^*(t)C(t)=0,\quad t\in\mathbb{R}.
			\end{split}
		\end{equation}
		admits a unique symmetric, positive semidefinite and $\theta$-periodic
		solution $P(\cdot)$, such that the associated closed-loop matrix  $A(\cdot)-B(\cdot)Q^{-1}(\cdot)B^*(\cdot)P(\cdot)$ is exponentially stable\footnote{Indeed, if $A(\cdot)$ is periodic and bounded, then $\dot y(t)=A(t)y(t)$ is exponentially stable if and only if for every fixed $s$, the norm of associated transition matrix $\|\Phi(t,s)\|$ approaches to zero as $t$ tends to infinity.},
		if and only if
		$(A(\cdot),B(\cdot))$ is exponentially $\theta$-periodic stabilizable
		and $(A(\cdot), C(\cdot))$ is exponentially $\theta$-periodic detectable.
	\end{lemma}

		\begin{equation}
			\begin{split}
				\dot{P}(t)&+A^*(t)P(t)+P(t)A(t)\\
				&-P(t)B(t)Q^{-1}(t)B^*(t)P(t)\\
				&+C^*(t)C(t)=0,\quad t\in\mathbb{R}.
			\end{split}
		\end{equation}

	\begin{remark}\label{du120505}
		Suppose that $(A(\cdot),B(\cdot))$ is exponentially $\theta$-periodic stabilizable
		and $(A(\cdot), C(\cdot))$ is exponentially $\theta$-periodic detectable.
		Let $P(\cdot)$ be the solution of \eqref{du12052} as guaranteed by Lemma \ref{ric}, and let $\Psi(\cdot,\cdot)$ denote the transition matrix (or evolution operator) associated with the feedback matrix $A(\cdot)-B(\cdot)Q^{-1}(\cdot)B^*(\cdot)P(\cdot)$. 
		Then, there exist positive constants $c$ and $\lambda$ such that the periodicity property
		\begin{equation}\label{du12051}
			\Psi(t+\theta,s+\theta)=\Psi(t,s) \;\;\;\text{for all}\;\;t>s
		\end{equation}
		and the exponential stability estimate
		\begin{equation}
			\|\Psi(t,s)\|\leq ce^{-\nu (t-s)}\;\;\;\text{for all}\;\;t>s
		\end{equation}
		hold true.
	\end{remark}

	\begin{lemma}\label{lya}
		Suppose that $(A(\cdot),B(\cdot))$ is exponentially $\theta$-periodic stabilizable
		and $(A(\cdot), C(\cdot))$ is exponentially $\theta$-periodic detectable. Let $P(\cdot)$ be the unique solution of \eqref{du12052}. Then, the periodic matrix Lyapunov differential equation
		\begin{multline*}
			-\dot{E}(t)+\big(A(t)-B(t)Q^{-1}(t)B^*(t)P(t)\big)E(t)\\
			+E(t)\big(A(t)-B(t)Q^{-1}(t)B^*(t)P(t)\big)^*\\
			-B(t)Q^{-1}(t)B^*(t)=0,\;\;t\in\mathbb R,
		\end{multline*}
		admits a unique symmetric, bounded, and $\theta$-periodic solution $E(\cdot)$, which is explicitly given by
		\begin{equation}\label{lya1}
			E(t)= -\int_{-\infty}^t\Psi(t,s)B(s)Q^{-1}(s)B^*(s)\Psi^*(t,s)\,\mathrm{d}s,
		\end{equation}
		where $\Psi(\cdot,\cdot)$ is the transition matrix associated with $A(\cdot)-B(\cdot)Q^{-1}(\cdot)B^*(\cdot)P(\cdot)$.
	\end{lemma}

	\begin{theorem}\label{dudic}
		Assume that $(A(\cdot),B(\cdot))$ is exponentially $\theta$-periodic stabilizable and that $(A(\cdot),C(\cdot))$ is exponentially $\theta$-periodic detectable. Then, the linear coupled system
		\begin{equation*}
			\begin{pmatrix}
				\dot{y}(t)\\\dot{\lambda}(t)
			\end{pmatrix}=
			\begin{pmatrix}
				A(t) & B(t)Q^{-1}(t)B^*(t)\\C^*(t)C(t)&-A^*(t)
			\end{pmatrix}\begin{pmatrix}
				y(t)\\ \lambda (t)
			\end{pmatrix},\; t\in\mathbb R,
		\end{equation*}
		can be decoupled into the block-diagonal form
		\begin{equation*}
			\begin{split}
				\begin{pmatrix}\dot{p}(t)\\ \dot{q}(t)\end{pmatrix}
				=
				\begin{pmatrix}
					L(t) & 0 \\
					0 & -L^*(t)
				\end{pmatrix}
				&
				\begin{pmatrix}p(t)\\ q(t)\end{pmatrix},\quad t\in\mathbb{R},
			\end{split}
		\end{equation*}
		where
		\begin{equation}\label{eq:A-BQ-1B*P}
			L(t)=A(t)-B(t)Q^{-1}(t)B^*(t)P(t),\quad t\in\mathbb{R}
		\end{equation}
		by using the linear, bounded, and $\theta$-periodic transformation
		\begin{equation}\label{ducan1}
			\begin{pmatrix}
				p(t)\\ q (t)
			\end{pmatrix}=
			\begin{pmatrix}
				I+E(t)P(t)&E(t)\\P(t)&I
			\end{pmatrix}
			\begin{pmatrix}
				y(t)\\  \lambda (t)
			\end{pmatrix},\quad t\in\mathbb R,
		\end{equation}
		where $P(\cdot)$ and $E(\cdot)$ are the solutions of the periodic Riccati and Lyapunov differential equations given in Lemmas \ref{ric} and \ref{lya}, respectively.
	\end{theorem}

		\begin{equation}
			\begin{pmatrix}
				p(t)\\ q (t)
			\end{pmatrix}=
			\begin{pmatrix}
				I+E(t)P(t)&E(t)\\P(t)&I
			\end{pmatrix}
			\begin{pmatrix}
				y(t)\\  \lambda (t)
			\end{pmatrix},\quad t\in\mathbb R,
		\end{equation}

	\begin{equation}\label{PG}
		\left\{
		\begin{aligned}
			&\dot{P}(t)+A^*(t)P(t)+P(t)A(t)-P(t)B(t)Q^{-1}(t)B^*(t)P(t)\\
			&\quad +C^*(t)C(t)=0,\quad t\in(-\infty,T],\\
			&P(T)=G>0.
		\end{aligned}
		\right.
	\end{equation}

	\begin{theorem}\label{thm:analy_riccati}
		Assume that $(A(\cdot),B(\cdot))$ is exponentially $\theta$-periodic stabilizable, and that $(A(\cdot),C(\cdot))$ is exponentially $\theta$-periodic detectable. Then, the solution $P(\cdot)$ of \eqref{PG} satisfies 
		\begin{equation}\label{eq:riccati_esti}
			\|P(t)-P_0(t)\|\leq C e^{-2\nu(T-t)}, \quad \forall t \in (-\infty, T],
		\end{equation}
		for some constants $C>0$ and $\nu>0$, where $P_0(\cdot)$ is the unique positive semidefinite $\theta$-periodic solution to the periodic Riccati equation \eqref{du12052}.
	\end{theorem}

	\begin{proof}[Proof of Theorem \ref{thm:analy_riccati}]
		Under the stabilizability and detectability assumptions, Lemma \ref{ric} guarantees the existence of a unique positive semidefinite $\theta$-periodic solution $P_0(\cdot)$ to \eqref{du12052}.
		
		To apply the periodic dichotomy transformation established in Theorem \ref{dudic}, we perform a change of variables to align the signs. Let $y(\cdot) = X(\cdot)$ and $\lambda(\cdot) = -Y(\cdot)$. The system governing $(X,Y)$ is then transformed into the standard Hamiltonian system:
		\begin{equation*}
			\begin{pmatrix}
				\dot{y}(t)\\ \dot{\lambda}(t)
			\end{pmatrix}=
			\begin{pmatrix}
				A(t) & B(t)Q^{-1}(t)B^*(t)\\
				C^*(t)C(t)&-A^*(t)
			\end{pmatrix}
			\begin{pmatrix}
				y(t)\\ \lambda (t)
			\end{pmatrix},  
		\end{equation*}for $t\in (-\infty,T]$.
		
		Applying the dichotomy transformation \eqref{ducan1}, we introduce the auxiliary variables
		\begin{equation*}
			\begin{aligned}
				\begin{pmatrix} \hat{X}(t) \\ \hat{Y}(t) \end{pmatrix} &:= 
				\begin{pmatrix}
					I+E(t)P_0(t) & E(t)\\
					P_0(t) & I
				\end{pmatrix}
				\begin{pmatrix} y(t) \\ \lambda(t) \end{pmatrix}\\
				&=
				\begin{pmatrix}
					I+E(t)P_0(t) & -E(t)\\
					P_0(t) & -I
				\end{pmatrix}
				\begin{pmatrix} X(t) \\ Y(t) \end{pmatrix}.
			\end{aligned}
		\end{equation*}
		This transformation decouples the system into
		\begin{equation*}
			\begin{pmatrix} \dot{\hat{X}}(t)\\ \dot{\hat{Y}}(t) \end{pmatrix}
			=
			\begin{pmatrix}
				L(t) & 0 \\
				0 & -L^*(t)
			\end{pmatrix}
			\begin{pmatrix} \hat{X}(t)\\ \hat{Y}(t) \end{pmatrix},\quad t\in(-\infty,T],
		\end{equation*}
		where $L(\cdot)=A(\cdot)-B(\cdot)Q^{-1}(\cdot)B^*(\cdot)P_0(\cdot)$ is the exponentially stable closed-loop matrix. Consequently,
		\begin{equation}
			\begin{aligned}
				\hat{X}(t) &= \Psi(t,T)\hat{X}(T), \\
				\hat{Y}(t) &= \Psi^*(T,t)\hat{Y}(T),
			\end{aligned}
		\end{equation}
		where $\Psi(\cdot,\cdot)$ is the transition matrix associated with $L(\cdot)$. 
		
		Using the inverse of the dichotomy transformation, we obtain
		\begin{equation}\label{inv_trans_XY}
			\begin{aligned}
				X(t) &= \hat{X}(t) - E(t)\hat{Y}(t), \\
				Y(t) &= P_0(t)\hat{X}(t) - \big(I+P_0(t)E(t)\big)\hat{Y}(t),
			\end{aligned}
		\end{equation} for $t\in(-\infty,T]$.
		Subtracting $P_0(\cdot)X(\cdot)$ from $Y(\cdot)$ yields the identity:
		\begin{equation}
			Y(t) - P_0(t)X(t) = -\hat{Y}(t),\quad \forall t\in(-\infty,T].
		\end{equation}
		Therefore,
		\begin{equation}\label{P_error_exact}
			P(t) - P_0(t) = Y(t)X^{-1}(t) - P_0(t) = -\hat{Y}(t)X^{-1}(t),
		\end{equation}for $t\in(-\infty,T]$.
		
		Next, we estimate the norm of the right-hand side. Substituting $\hat{X}(\cdot) = \Psi(\cdot,T)\hat{X}(T)$ and $\hat{Y}(\cdot) = \Psi^*(T,\cdot)\hat{Y}(T)$ into the expression for $X(\cdot)$ yields
		\begin{equation}
			X(t) = \Psi(t,T)\hat{X}(T) - E(t)\Psi^*(T,t)\hat{Y}(T),\quad\forall t\in(-\infty,T].
		\end{equation}
		For $t\in(-\infty,T]$, define
		\begin{equation}
			S(t) := \Psi(T,t)X(t) = \hat{X}(T) - \Psi(T,t)E(t)\Psi^*(T,t)\hat{Y}(T).
		\end{equation}
		From Remark \ref{du120505}, we have the exponential decay estimate $\|\Psi(T,t)\| \leq c e^{-\nu(T-t)}$. Since $E(\cdot)$ is bounded, the term $\Psi(T,t)E(t)\Psi^*(T,t)\hat{Y}(T)$ decays exponentially to zero as $t \to -\infty$. Assuming the terminal matrix $\hat{X}(T) = I + E(T)(P_0(T)-G)$ is invertible (which holds generically), it follows that $S(t)$ is invertible for $T-t$ sufficiently large, and its inverse $S^{-1}(t)$ is uniformly bounded. 
		
		Since $X(\cdot) = \Psi(\cdot,T)S(\cdot)$, we can write $X^{-1}(\cdot) = S^{-1}(\cdot)\Psi(T,\cdot)$. Substituting this and $\hat{Y}(\cdot) = \Psi^*(T,\cdot)\hat{Y}(T)$ into \eqref{P_error_exact}, we obtain
		\begin{equation}
			P(t) - P_0(t) = -\Psi^*(T,t)\hat{Y}(T) S^{-1}(t) \Psi(T,t),
		\end{equation}for $t\in(-\infty,T]$.
		Taking the norm on both sides, we deduce
		\begin{equation}
			\begin{split}
				\|P(t) - P_0(t)\| &\leq \|\Psi^*(T,t)\| \|\hat{Y}(T)\| \|S^{-1}(t)\| \|\Psi(T,t)\| \\
				&\leq c^2 \|\hat{Y}(T)\| \|S^{-1}(t)\| e^{-2\nu(T-t)}.
			\end{split}
		\end{equation}
		Since $\|S^{-1}(t)\|$ is uniformly bounded, there exists a constant $C > 0$ such that
		\begin{equation}
			\|P(t) - P_0(t)\| \leq C e^{-2\nu(T-t)}, \quad \forall t \in (-\infty, T].
		\end{equation}
		This completes the proof.
	\end{proof}

\end{document}
